\documentclass[11pt,a4paper]{article}

\usepackage[margin=2.5cm]{geometry}
\usepackage{amsmath,amssymb,bm}
\usepackage{xcolor}
\usepackage[colorlinks,linkcolor=blue,urlcolor=blue]{hyperref}

\newcommand{\vek}[1]{\bm{#1}}
\newcommand{\grad}{\vek\nabla}
\newcommand{\impl}{\textcolor{green!60!black}{Implemented.}}
\newcommand{\referee}[1]{\medskip\noindent{\sffamily\large #1)~Referee:}\par\smallskip}
\newcommand{\reply}{\par\medskip\noindent{\sffamily\large Reply:}\par\smallskip\noindent}
\newenvironment{reftext}{\itshape}{\par}

\begin{document}

\begin{center}
{\sffamily\LARGE Reply to referee report on manuscript EPJC-26-08-339}
\end{center}

\bigskip

\noindent We would like to thank the referee for taking the time to read our manuscript carefully, for the positive overall evaluation, and for constructive suggestions that help us further improve the presentation. We give a detailed response to all the points raised by the referee below. For the convenience of the editor and the referee, we attach a marked-up version of the manuscript, where modified parts of the original text are highlighted in color. With these changes, we hope that the revised manuscript can be accepted for publication.

%-------------------------------------------------------------------------------
\referee{1}
\begin{reftext}
Isn't the notation $H$ unfortunate because it usually refers to the auxiliary magnetic field in the presence of a medium? Although there is a medium here I think $H$ does not have this usual meaning, which may cause a slight, but unnecessary, confusion.
\end{reftext}

\reply
We understand how this notation may lead to confusion. In the revised manuscript, we use the symbol $\vek b$ for the rescaled magnetic field instead of $\vek H$, which hopefully does not collide with any other generally accepted notation. In case the referee feels that there is a better choice of symbol, we kindly ask them to let us know; the notation can be changed at once in the whole paper by redefining a single \LaTeX{} command in the preamble of the manuscript. \impl

%-------------------------------------------------------------------------------
\referee{2}
\begin{reftext}
I am confused by the discussion at the end of Sec 2 and by Fig 1. Eqs (14) and (15) give two different classes of solutions with $\phi$ nonzero. Then there is the third possibility of $\phi=0$. Why is the transition between the former two identified with the analog of the CSL transition? Does that mean in Fig 1 there is a $\phi\neq0$ state everywhere, below and above the curve? I suspect I am misunderstanding the explanation, but since other readers might have the same problem, a clarification would be useful.
\end{reftext}

\reply
In a finite volume, the ground state will always have nonzero $\phi$, as required by the natural boundary condition~(16). The trivial vacuum $\phi_0=0$ does not satisfy this boundary condition and thus is not even a stationary point of the energy functional. As a consequence, the ground state will have negative total energy for any (nonzero) magnetic field and any choice of the system size. We therefore need a different criterion to locate the phase transition than a simple comparison of the total energy to that of the trivial vacuum state. In order to help the reader understand the content of the phase diagram, we have added a new figure, which compares profiles of the ground state just below and just above the phase transition. This explains visually the difference between the ``vacuum solution'' (mostly constant $\phi_0$, bent near the boundary so as to satisfy the natural boundary condition) and the ``domain-wall-like CSL solution'' (maximum gradient at the origin). We have also modified the discussion at the end of Sec.~2.2 so as to accommodate the new figure. \impl

%-------------------------------------------------------------------------------
\referee{3}
\begin{reftext}
Are there any physical lessons to be learned from the obtained configurations, say in Figs 3 and 4? I understand that the focus of the paper is mostly on developing the mathematical tools, but the reader might wonder if the behavior of the baryon number is expected or surprising etc. Some physical explanation/intuition would then also be useful to predict the qualitative behavior of configurations with more complicated inhomogeneities in the magnetic field.
\end{reftext}

\reply
It is indeed the case that our choice of the magnetic field with a Gaussian profile (Fig.~4 in the revised manuscript) was meant just for illustration. Similar, but somewhat different, localized profiles are sometimes used to model magnetic fields present in heavy-ion collisions. Based on our numerical experience, we believe that drawing general lessons from a single solution is subtle. Namely, it is clear from Fig.~4 that the CSL solution is highly sensitive to boundary conditions: the four-fold symmetry of the baryon number density is not present in the Ansatz for the magnetic field, and stems from minimizing the energy functional on a square-shaped domain. The relation between the magnetic field and the CSL solution is therefore, generally speaking, highly nonlocal. (In the chiral limit, the CSL solution is entirely fixed by the values of the magnetic field on the boundary.) The only generic exception is the case of conservative magnetic fields in the chiral limit where, as explained in Sec.~3.1, the CSL state $\phi_0$ minimizes the energy functional pointwise, and the gradient $\grad\phi_0$ equals $\vek B$ everywhere up to a common conversion factor. We have added some comments in this regard to the part of the text referring to Fig.~4. \impl

The status of Fig.~5 (in the new numbering) is somewhat different. This was meant to complement the analytical computation of the CSL solution in the presence of a magnetic domain wall. In the chiral limit, a uniform magnetic field leads to a CSL ground state with constant gradient; cf.\ Eq.~(4). The figure highlights the deformation of this solution due to the presence of a localized domain wall. In case the referee thinks that the figure does not convey any insight beyond the analytic solution in Eq.~(24), we would be happy to simply remove the figure. Please note that in the revised manuscript, we have made the caption of Fig.~5 more precise by giving an explanation of the shading, which is somewhat different than in Fig.~4. Namely, for the domain wall, the baryon number density is nonnegative everywhere. It vanishes along the domain wall and reaches its largest value, corresponding to that in a uniform magnetic field, only on the boundaries $z=\pm L_z/2$. In Fig.~5, dark shading thus indicates vanishing baryon number density and light shading its largest positive values, whereas in Fig.~4, the dark regions represent negative baryon number density. \impl

%-------------------------------------------------------------------------------
\referee{4}
\begin{reftext}
The authors mention in the introduction that inhomogeneities in $B$ are present in heavy-ion collisions and in neutron stars but that at least in the latter variations are on a scale much larger than relevant for modification of the original CSL state. It might be helpful to give some numbers for the scales involved to get an idea how far (if at all) the configurations seen here are away from what we would expect in physical systems. Also, would these estimates change if neutron star mergers are considered, where there are magneto-rotational instabilities and thus small-scale variations of the magnetic field? Presumably even these scales are much larger than what is needed for the configurations considered here, and also temperature effects in mergers would destroy any CSL-type phase?
\end{reftext}

\reply
We have added a paragraph with order-of-magnitude estimates to the introduction. The scale of the gradient of the CSL solution is set by the Compton wavelength of the pion, $1/m_\pi\approx1.4~\text{fm}$. For the inhomogeneity of the magnetic field to be of relevance to the CSL state, it should vary significantly at length scales of the same order of magnitude or shorter. In heavy-ion collisions, the magnetic field does vary over a few fm, but it decays on time scales of the order of fm$/c$, the baryon chemical potential is small and the temperature is too high for the CSL to survive. In neutron stars, the magnetic field is likely not strong enough to support a CSL ground state, and its scale of nonuniformity is set by the size of the star, that is, is of the order of km. In neutron star mergers, the Kelvin--Helmholtz and magnetorotational instabilities amplify the field to $\gtrsim10^{16}~\text{G}$ and generate small-scale turbulence. However, ``small-scale'' here means small compared to the size of the star, not compared to the scale of QCD. For instance, the wavelength of the fastest-growing mode of the magnetorotational instability, $\lambda_\text{MRI}\simeq2\pi v_\text{A}/\Omega$, where $v_\text{A}=B/\sqrt{4\pi\rho}$ is the Alfv\'en velocity and $\Omega$ the angular velocity, ranges from about a meter for $B\sim10^{13}~\text{G}$ to kilometers for $B\gtrsim10^{15}~\text{G}$, for densities $\rho\sim10^{13}$--$10^{14}~\text{g}/\text{cm}^3$ and angular velocities $\Omega\sim10^3$--$10^4~\text{s}^{-1}$ typical of merger remnants~\cite{Siegel2013MRI,Guilet2017MRI}. Likewise, the Kelvin--Helmholtz instability grows fastest at the shortest scales, which in numerical simulations are set by the grid resolution of order $10~\text{m}$~\cite{Kiuchi:2015sga}. These scales exceed $1/m_\pi$ by at least fifteen orders of magnitude. Moreover, the temperatures of several tens of MeV reached in merger remnants would likely destroy any CSL-type order. We therefore do not believe that the sample magnetic fields studied in our paper are of direct relevance to the phenomenology of neutron stars and their mergers. The most likely domain of applications would be lattice QCD simulations with nonuniform magnetic fields. We have added the above estimate of the MRI wavelength to the introduction. \impl

%-------------------------------------------------------------------------------
\referee{5}
\begin{reftext}
Why does the EoM for $\phi$ not have to be coupled with Maxwell's equations? I understand that the coupling of the pions to $B$ is anomalous, so presumably Maxwell's equations are trivially fulfilled? Even if that's the case maybe a short remark would be helpful.
\end{reftext}

\reply
This paper assumes, in line with most literature on the CSL phase in quark matter, that the magnetic field is fixed and does not experience any feedback from the neutral pion condensate. The source of the magnetic field is not included in the Lagrangian and may constitute other type of matter than pions. One therefore cannot expect the magnetic field to satisfy the Maxwell equations that would follow by adding to our Lagrangian the kinetic term for the electromagnetic field. We have added a remark in this regard to the conclusions. \impl{} Here, we would like to note that it is possible to consider a dynamical coupling between the electromagnetic field and pions, and it leads to interesting and highly nontrivial effects well beyond the scope of the present paper. An early reference in this regard is Brauner and Kadam, JHEP \textbf{03} (2017) 015, which showed that, in analogy with axion electrodynamics, the anomalous term in the Lagrangian leads to a mixing between neutral pions and photons in the presence of a background magnetic field, which has dramatic consequences for the photon spectrum as well as for the low-temperature thermodynamics of the system.

%-------------------------------------------------------------------------------
\begin{thebibliography}{9}\small
\bibitem{Siegel2013MRI} D.~M. Siegel, R.~Ciolfi, A.~I. Harte, L.~Rezzolla, Phys.\ Rev.\ D \textbf{87} (2013) 121302, arXiv:1302.4368.
\bibitem{Guilet2017MRI} J.~Guilet, A.~Bauswein, O.~Just, H.-T. Janka, Mon.\ Not.\ Roy.\ Astron.\ Soc.\ \textbf{471} (2017) 1879, arXiv:1610.08532.
\bibitem{Kiuchi:2015sga} K.~Kiuchi, P.~Cerd\'a-Dur\'an, K.~Kyutoku, Y.~Sekiguchi, M.~Shibata, Phys.\ Rev.\ D \textbf{92} (2015) 124034, arXiv:1509.09205.
\end{thebibliography}

\end{document}
